%%%%%%%%%%%%%%%%%%%%%%%%%%%%%%%%%%%%%%%%%
%%%  Section:
%%%%%%%%%%%%%%%%%%%%%%%%%%%%%%%%%%%%%%%%%
\section{Selected Research Highlights: Quantitative Impact and Results}

\begin{comment}
\subsection{High-Impact Breakthrough: Overcoming the UQ Bottleneck \\
	Accelerating Uncertainty Quantification in Nonlinear Circuits Simulation}

%%%%%%%%%%%%%%%%%%%%%%
%%%     Note       %%%
%%%%%%%%%%%%%%%%%%%%%%
\begin{frame}<0>{NOTES TO PRESENTER:}
	"In the world of 5G and high-speed design, 'worst-case corners' aren't enough—we need Monte Carlo. But as we all know, Newton-Raphson iterations at every sample point kill performance. My research on parameterized DEIM (Slide 10) attacks this specific bottleneck. By achieving a 50x speedup, we allow engineers to actually run these checks before tape-out, potentially saving millions in silicon re-spins."
\end{frame}
\end{comment}

\begin{comment}
%%%%%%%%%%%%%%%%%%%%%%
%%%     Note       %%%
%%%%%%%%%%%%%%%%%%%%%%
\begin{frame}<0>{NOTES TO PRESENTER:}
	"In the world of 5G and high-speed design, 'worst-case corners' aren't enough—we need Monte Carlo. But as we all know, Newton-Raphson iterations at every sample point kill performance. My research on parameterized DEIM (Slide 10) attacks this specific bottleneck. By achieving a 50x speedup, we allow engineers to actually run these checks before tape-out, potentially saving millions in silicon re-spins."While my previous slide discussed DC-centric bottlenecks, our work in the Frequency Domain is equally transformative. In Harmonic Balance, the 'Transformation Tax'—constantly switching between time and frequency—is what prevents us from simulating large-scale RFICs. By extending DEIM directly into the HB framework (Slide 11), we achieved a 30x runtime reduction. For a design team, this is the difference between waiting a week for results or having them before the afternoon coffee break."
\end{frame}
\end{comment}

\begin{comment}
%%%%%%%%%%%%%%%%%%%%%%
%%%     Note       %%%
%%%%%%%%%%%%%%%%%%%%%%
\begin{frame}<0>{NOTES TO PRESENTER:}
	"In my final technical highlight (Slide 12), I want to address the most critical bottleneck for any EDA tool: Stability. We can make a solver 100x faster, but if it crashes due to an unstable reduced-order model, it will never be used in production. My research ensures Stability by Construction. By solving the stability constraint as a generalized eigenproblem rather than a slow optimization task, we provide a robust, sparse, and fail-safe model for active circuits. This is the level of mathematical integrity I intend to bring to the Siemens EDA Chair."

	Why this completes your Technical Narrative:
	Slide 10 (UQ): Proves you can handle Statistical Complexity.

	Slide 11 (HB): Proves you can handle Domain Complexity (RF/Periodic).

	Slide 12 (Stability): Proves you can handle Numerical Integrity.
\end{frame}

%%%%%%%%%%%%%%%%%%%%%%
%%%     Note       %%%
%%%%%%%%%%%%%%%%%%%%%%
\begin{frame}<0>{NOTES TO PRESENTER:}
	"In this final highlight (Slide 13), we look at the intersection of Electromagnetics and Circuits. PEEC modeling is essential for high-frequency design, but the propagation delays turn it into a nightmare of Neutral Delay Differential Equations. My research replaces slow, unstable time-stepping with a compact PMOR and NILT-based solver (Slide 13). The result? We achieved a three-order-of-magnitude speedup. This effectively moves large-scale retarded systems from the 'prohibitively slow' category into the 'daily design task' category."

	Summary of your 4-Slide Technical "Punch":
	Slide 10 (DC/UQ): 50x speedup for Statistical Analysis.

	Slide 11 (HB/Periodic): 30x speedup for RFIC Analysis.

	Slide 12 (Stability): Guaranteed Integrity for Active Circuits.

	Slide 13 (PEEC/Delay): 1000x speedup for EM/Retarded Systems.
\end{frame}
\end{comment}

\begin{comment}
%%%%%%%%%%%%%%%%%%%%%%
%%%     Slide      %%%
%%%%%%%%%%%%%%%%%%%%%%
\begin{frame}{EM-Circuit Co-Simulation: Delayed PEEC Systems -- PMOR and NILT for High-Frequency Retarded Networks}

	\lmv{-1}

	\begin{mybox}
		\nlt
		B.~Nouri, L.~Lombardi, Y.~Tao, M.~Nakhla, F.~Ferranti, and G.~Antonini,
		``Parameterized model order reduction of delayed PEEC circuits,''
		\emph{IEEE Trans. Electromagn. Compat.},
		vol.~62, no.~3, pp.~859--869, Jun.~2020.
	\end{mybox}

	\vspace{0.3cm}

	\begin{columns}[T]

		\begin{column}{0.65\textwidth}
			\small
			\textbf{The Challenge: Neutral Delay DEs (NDDEs)}
			\begin{itemize}
				\item High-speed interconnects require PEEC with retardation.
				\item Time-stepping is bottlenecked by constant interpolation.
				\item Instability and passivity issues are common in NDDEs.
			\end{itemize}

			\vspace{0.2cm}

			\textbf{The Innovation: Decoupled PMOR + NILT}
			\begin{itemize}
				\item \textbf{Compact Subspaces:}
				      Frequency-independent moment matching for minimal ROM size.
				\item \textbf{NILT Solver:}
				      Response evaluation in the Laplace domain, bypassing time-domain interpolation entirely.
			\end{itemize}
		\end{column}

		\begin{column}{0.3\textwidth}
			\begin{alertblock}{{\Large\green{Measured Impact:}}}
				\begin{mybox}[red]
					\centering
					\vspace{0.1cm}
					{\huge \textbf{1000$\times$}}\\
					\small \textbf{Simulation Speedup}\\[0.4cm]
					{\large \textbf{NILT vs.\ NDDE}}\\
					\small \textbf{Solver Robustness}
				\end{mybox}
			\end{alertblock}
		\end{column}

	\end{columns}

	\vfill
	\begin{center}
		\red{\bfseries\large
			``Scaling the simulation of retarded systems by three orders of magnitude
			for high-frequency SI/PI verification.''}
	\end{center}
\end{frame}
\end{comment}

% %%%%%%%%%%%%%%%%%%%%%%
% %%%     Slide      %%%
% %%%%%%%%%%%%%%%%%%%%%%
% \begin{frame}[t, fragile]{\textbf{UQ:~} Efficient DC-Centric Simulation of Nonlinear Systems Under Uncertainty}
% 	\lmv{-1}
% 	\begin{mybox}
% 		\nlt B.~Nouri, E.~Gad, A.~Nouri, and M.~Nakhla, ``{DC}-centric parameterized reduced-order model via moment-based interpolation projection {(MIP)} algorithm,'' \emph{{IEEE} Trans. Compon., Packag., Manuf. Technol.}, vol.~10, no.~8, pp.~1348--1357, Aug.~2020.
% 	\end{mybox}
% 	\vspace{0.4cm}
% 	\begin{columns}[T]
% 		\begin{column}{0.65\textwidth}

% 			\textbf{Challenge: The Newton-Raphson Bottleneck}
% 			\begin{itemize}
% 				\item Standard MC is too slow for SI/PI \& RF.
% 				\item Iterative solvers recompute DC points for every sample.
% 				\item Evaluation of reduced nonlinear terms is the primary cost.
% 			\end{itemize}

% 			\vspace{0.2cm}
% 			\textbf{Solution: Parameterized DEIM Formulation}
% 			\begin{itemize}
% 				\item Developed a surrogate using \textbf{DC-centric DEIM}.
% 				\item Directly exploits structurally parameterized moments.
% 			\end{itemize}
% 		\end{column}

% 		\begin{column}{0.3\textwidth}
% 			\begin{alertblock}{\green{\bfseries\Large~~~Measured Impact:}}
% 				\begin{mybox}[red]
% 					\centering
% 					\vspace{0.1cm}
% 					\noindent{\huge\textbf{Up to 50$\times$}}\par\smallskip
% 					\noindent{\small\textbf{Simulation Speed-up}}\par
% 					\vspace{0.4cm}
% 					\noindent {\large\textbf{Zero-Risk}} \par\smallskip
% 					\noindent \small \textbf{Tape-out Readiness}
% 				\end{mybox}
% 			\end{alertblock}
% 			\vspace{0.1cm}
% 		\end{column}
% 	\end{columns}

% 	\vfill
% 	\begin{center}
% 		\red{\bfseries \large %
% 			"Moving UQ from a 'prohibitive cost' to a 'practical design-stage tool' for SI/PI \& RF systems."
% 		}
% 	\end{center}
% \end{frame}

% %%%%%%%%%%%%%%%%%%%%%%
% %%%     Slide      %%%
% %%%%%%%%%%%%%%%%%%%%%%
% \begin{frame}{\textbf{HB:~} RFIC Innovation: Efficient Periodic Steady-State (PSS) Analysis}
% 	\lmv{-1}
% 	\begin{mybox}
% 		\nlt B.~Nouri, E.~Gad, A.~Nouri, and M.~Nakhla, ``Parameterized periodic steady-state analysis via reduced-order modelling,'' \emph{{IEEE} Trans. Microw. Theory Tech.}, vol.~69, no.~11, pp.~4825--4836, Nov.~2021.
% 	\end{mybox}

% 	\vspace{0.3cm}

% 	\begin{columns}[T]
% 		\begin{column}{0.6\textwidth}
% 			\small
% 			\textbf{Problem: The "Transformation Tax"}
% 			\begin{itemize}
% 				\item HB analysis is essential for RF/Microwave steady-state.
% 				\item Traditional solvers "ping-pong" between Time and Frequency domains for nonlinear terms.
% 				\item FFT/IFFT operations dominate the CPU time.
% 			\end{itemize}

% 			\vspace{0.2cm}
% 			\textbf{Innovation: Direct Harmonic Domain DEIM}
% 			\begin{itemize}
% 				\item Evaluation of surrogates directly in the harmonic domain.
% 				\item Eliminates domain conversion overhead entirely.
% 				\item Successfully benchmarked on large-scale RF designs.
% 			\end{itemize}
% 		\end{column}

% 		\begin{column}{0.3\textwidth}
% 			\begin{alertblock}{\Large\green{~~~Measured Impact:}}
% 				\begin{mybox}[red]
% 					\centering				\centering
% 					\vspace{0.1cm}
% 					{\huge \textbf{Up to 30$\times$}} \\
% 					\small \textbf{Total Speedup} \\
% 					\vspace{0.4cm}
% 					{\large \textbf{Days $\rightarrow$ Hours}} \\
% 					\small \textbf{Design Cycle Time}
% 				\end{mybox}
% 			\end{alertblock}
% 			\vspace{0.2cm}
% 		\end{column}
% 	\end{columns}
% 	\vfill
% 	\begin{center}
% 		\red{\bfseries\large %
% 			"Transforming infeasible multi-day simulations into routine, automated verification tasks for high-yield RFICs."}
% 	\end{center}
% \end{frame}

% % %%%%%%%%%%%%%%%%%%%%%%
% % %%%     Slide      %%%
% % %%%%%%%%%%%%%%%%%%%%%%
% % \begin{frame}{EM-Circuit Co-Simulation: Delayed PEEC Systems -- PMOR and NILT for High-Frequency Retarded Networks}

% % \lmv{-1}

% % \begin{mybox}
% % 	\nlt
% % 	B.~Nouri, L.~Lombardi, Y.~Tao, M.~Nakhla, F.~Ferranti, and G.~Antonini,
% % 	``Parameterized model order reduction of delayed PEEC circuits,''
% % 	\emph{IEEE Trans. Electromagn. Compat.},
% % 	vol.~62, no.~3, pp.~859--869, Jun.~2020.
% % \end{mybox}

% %%%%%%%%%%%%%%%%%%%%%%
% %%%     Slide      %%%
% %%%%%%%%%%%%%%%%%%%%%%
% %%%%%%%%%%%%%%%%%%%%%%%%%%%%%%%%%%%%%%%%%%%%%%%%%%%%%%%%%%%%%%%%%%
% %%% Slide: Stability by Construction
% %%%%%%%%%%%%%%%%%%%%%%%%%%%%%%%%%%%%%%%%%%%%%%%%%%%%%%%%%%%%%%%%%%
% \begin{frame}[t]{Stability by Construction: Active Circuit Macromodeling}
% \lmv{-1}
% \begin{mybox}
% 	\scriptsize
% 	G. Ghaly, E. Gad, M. Nakhla, and \textbf{B. Nouri}, ``An Algorithmic Approach to Formulate Well-Conditioned Stable Reduced-Order Models of Active Circuits,'' \textit{IEEE Trans. CPMT}.
% \end{mybox}

% \vspace{1ex}

% \begin{columns}[T, onlytextwidth]

% %% LEFT COLUMN: The Narrative
% \begin{column}{0.62\textwidth}
% 	\small
% 	\textbf{\textcolor{blue!80!black}{The Problem: Simulation Divergence}}
% 	\begin{itemize}
% 		\item Active circuits frequently yield unstable ROMs under standard Galerkin projection.
% 		\item Results in catastrophic simulation failure during transient or sensitivity analysis.
% 	\end{itemize}

% 	\vspace{2.5ex}

% 	\textbf{\textcolor{blue!80!black}{The Innovation: Lyapunov-Based Projection}}
% 	\begin{itemize}
% 		\item Left-projection matrix ($V$) specifically constructed as a function of $W$ to satisfy \textbf{Lyapunov stability}.
% 		\item Formulated as an efficient \textbf{generalized eigenproblem}; avoids costly non-linear optimization.
% 		\item Native ill-conditioning correction for extreme numerical robustness.
% 	\end{itemize}
% \end{column}

% %% RIGHT COLUMN: The Result
% \begin{column}{0.34\textwidth}
% \setbeamercolor{block title example}{fg=white, bg=red!20%%%%%%%%%%%%%%%%%%%%%%%%%%%%%%%%%%%%%%%%%%%%%%%%%%%%%%%%%%%%%%%%%%
% %%% Slide: Stability by Construction
% %%%%%%%%%%%%%%%%%%%%%%%%%%%%%%%%%%%%%%%%%%%%%%%%%%%%%%%%%%%%%%%%%%
% \begin{frame}[t]{Stability by Construction: Active Circuit Macromodeling}
% 	\vspace{1ex}

% 	% Clean Citation Block
% 	\setbeamercolor{block body}{bg=gray!10, fg=black}
% 	\begin{block}{}
% 		\scriptsize
% 		G. Ghaly, E. Gad, M. Nakhla, and \textbf{B. Nouri}, ``An Algorithmic Approach to Formulate Well-Conditioned Stable Reduced-Order Models of Active Circuits,'' \textit{IEEE Trans. CPMT}.
% 	\end{block}

% 	\vspace{1ex}

% 	\begin{columns}[T, onlytextwidth]

% 		%% LEFT COLUMN: The Narrative
% 		\begin{column}{0.62\textwidth}
% 			\small
% 			\textbf{\textcolor{blue!80!black}{The Problem: Simulation Divergence}}
% 			\begin{itemize}
% 				\item Active circuits frequently yield unstable ROMs under standard Galerkin projection.
% 				\item Results in catastrophic simulation failure during transient or sensitivity analysis.
% 			\end{itemize}

% 			\vspace{2.5ex}

% 			\textbf{\textcolor{blue!80!black}{The Innovation: Lyapunov-Based Projection}}
% 			\begin{itemize}
% 				\item Left-projection matrix ($V$) specifically constructed as a function of $W$ to satisfy \textbf{Lyapunov stability}.
% 				\item Formulated as an efficient \textbf{generalized eigenproblem}; avoids costly non-linear optimization.
% 				\item Native ill-conditioning correction for extreme numerical robustness.
% 			\end{itemize}
% 		\end{column}

% 		%% RIGHT COLUMN: The Result
% 		\begin{column}{0.34\textwidth}
% 			\setbeamercolor{block title example}{fg=white, bg=green!45!black}
% 			\setbeamercolor{block body example}{bg=green!5, fg=black}
% 			\begin{exampleblock}{\centering \scriptsize \textbf{KEY BREAKTHROUGHS}}
% 				\centering
% 				\vspace{1ex}
% 				{\huge \textbf{60$\times$}} \\
% 				\footnotesize \textbf{Total Speedup} \\
% 				\vspace{1.5ex}
% 				\textcolor{green!30}{\rule{0.8\textwidth}{0.5pt}} \\
% 				\vspace{1.5ex}
% 				\scriptsize
% 				\textbf{Guaranteed Preservation:} \\
% 				\textcolor{blue!80!black}{\textbf{Asymptotic Stability}} \\
% 				\textit{(By Construction)}

% 				\vspace{2ex}

% 				\textcolor{blue!80!black}{\textbf{Numerical Robustness}} \\
% 				\textit{Well-Conditioned ROMs}
% 				\vspace{1ex}
% 			\end{exampleblock}
% 		\end{column}
% 	\end{columns}

% 	\vfill
% 	\vfill	%% Visionary Footer
% 	\begin{beamercolorbox}[wd=\textwidth, sep=1.2ex, center, rounded=true]{block body example}
% 		\small \itshape \textcolor{blue!80!black}{"Moving beyond 'fast' simulation to 'guaranteed' simulation for large-scale active networks."}
% 	\end{beamercolorbox}
% 	\vspace{1ex}
% 	%%%%%%%%%%%%%%%%%%%%%%%%%%%%%%%%%%%%%%%%%%%%%%%%%%%%%%%%%%%%%%%%%%
% 	%%% Slide: Stability by Construction
% 	%%%%%%%%%%%%%%%%%%%%%%%%%%%%%%%%%%%%%%%%%%%%%%%%%%%%%%%%%%%%%%%%%%
% 	\begin{frame}[t]{Stability by Construction: Active Circuit Macromodeling}
% 		\vspace{1ex}

% 		% Clean Citation Block
% 		\setbeamercolor{block body}{bg=gray!10, fg=black}
% 		\begin{block}{}
% 			\scriptsize
% 			G. Ghaly, E. Gad, M. Nakhla, and \textbf{B. Nouri}, ``An Algorithmic Approach to Formulate Well-Conditioned Stable Reduced-Order Models of Active Circuits,'' \textit{IEEE Trans. CPMT}.
% 		\end{block}

% 		\vspace{1ex}

% 		\begin{columns}[T, onlytextwidth]

% 			%% LEFT COLUMN: The Narrative
% 			\begin{column}{0.62\textwidth}
% 				\small
% 				\textbf{\textcolor{blue!80!black}{The Problem: Simulation Divergence}}
% 				\begin{itemize}
% 					\item Active circuits frequently yield unstable ROMs under standard Galerkin projection.
% 					\item Results in catastrophic simulation failure during transient or sensitivity analysis.
% 				\end{itemize}

% 				\vspace{2.5ex}

% 				\textbf{\textcolor{blue!80!black}{The Innovation: Lyapunov-Based Projection}}
% 				\begin{itemize}
% 					\item Left-projection matrix ($V$) specifically constructed as a function of $W$ to satisfy \textbf{Lyapunov stability}.
% 					\item Formulated as an efficient \textbf{generalized eigenproblem}; avoids costly non-linear optimization.
% 					\item Native ill-conditioning correction for extreme numerical robustness.
% 				\end{itemize}
% 			\end{column}

% 			%% RIGHT COLUMN: The Result
% 			\begin{column}{0.34\textwidth}
% 				\setbeamercolor{block title example}{fg=white, bg=green!45!black}
% 				\setbeamercolor{block body example}{bg=green!5, fg=black}
% 				\begin{exampleblock}{\centering \scriptsize \textbf{KEY BREAKTHROUGHS}}
% 					\centering
% 					\vspace{1ex}
% 					{\huge \textbf{60$\times$}} \\
% 					\footnotesize \textbf{Total Speedup} \\
% 					\vspace{1.5ex}
% 					\textcolor{green!30}{\rule{0.8\textwidth}{0.5pt}} \\
% 					\vspace{1.5ex}
% 					\scriptsize
% 					\textbf{Guaranteed Preservation:} \\
% 					\textcolor{blue!80!black}{\textbf{Asymptotic Stability}} \\
% 					\textit{(By Construction)}

% 					\vspace{2ex}

% 					\textcolor{blue!80!black}{\textbf{Numerical Robustness}} \\
% 					\textit{Well-Conditioned ROMs}
% 					\vspace{1ex}
% 				\end{exampleblock}
% 			\end{column}
% 		\end{columns}

% 		\vfill

% 		%% Visionary Footer
% 		\begin{beamercolorbox}[wd=\textwidth, sep=1.2ex, center, rounded=true]{block body example}
% 			\small \itshape \textcolor{blue!80!black}{"Moving beyond 'fast' simulation to 'guaranteed' simulation for large-scale active networks."}
% 		\end{beamercolorbox}
% 		\vspace{1ex}
% 	\end{frame}%%%%%%%%%%%%%%%%%%%%%%%%%%%%%%%%%%%%%%%%%
%%% Section: Distinguished Recognition
%%%%%%%%%%%%%%%%%%%%%%%%%%%%%%%%%%%%%%%%%
\section{Distinguished Research \& Recognized Excellence}

%%%%%%%%%%%%%%%%%%%%%%%%%%
%%%  Slide:
%%%%%%%%%%%%%%%%%%%%%%%%%%
\begin{frame}[t]{\insertsectionhead}
	\vspace{1ex}

	%% Main Block: Prestigious Awards & Invited Works
	\setbeamercolor{block title example}{fg=white, bg=blue!80!black}
	\begin{exampleblock}{\small \textbf{TOP-TIER HONORS \& AWARDS}}
		\vspace{1ex}
		\small
		\begin{columns}[t, onlytextwidth]
			\begin{column}{0.48\textwidth}
				\textbf{\textcolor{blue!80!black}{Best Transactions Paper (2013)}} \\
				\textit{IEEE Trans. CPMT} \\
				\textbf{Legacy:} Pioneered geometric MOR stopping criteria; now an industry standard.

				\vspace{1.5ex}

				\textbf{\textcolor{blue!80!black}{Invited State-of-the-Art (2021)}} \\
				\textit{IEEE Journal of Microwaves} \\
				\textbf{Impact:} Defined the global R\&D roadmap for next-gen EDA tools.
			\end{column}

			\begin{column}{0.48\textwidth}
				\textbf{\textcolor{blue!80!black}{Best Paper Award (2018)}} \\
				\textit{IEEE SPI, Brest} \\
				Time-domain sensitivity analysis.

				\vspace{1.5ex}

				\textbf{\textcolor{blue!80!black}{Best Paper Award (2011)}} \\
				\textit{IEEE SPI, Naples} \\
				Projection-based MOR stopping criteria.
			\end{column}
		\end{columns}
		\vspace{1ex}
	\end{exampleblock}

	%% Secondary Section: Permanent Literature & Stability Research
	\begin{columns}[t, onlytextwidth]

		% Left Side: Permanent Literature
		\begin{column}{0.48\textwidth}
			\setbeamercolor{block title}{fg=white, bg=green!45!black}
			\begin{block}{\small \textbf{PERMANENT LITERATURE}}
				\vspace{0.5ex}
				\small
				\textbf{\textcolor{blue!80!black}{Handbook Chapter (2020)}} \\
				\textit{"MOR in Microelectronics"} \\
				\textbf{In:} \textit{Handbook on MOR}, De Gruyter. \\
				\textcolor{gray}{Foundational reference for the EDA community.}
			\end{block}
		\end{column}

		% Right Side: Stability Pillars
		\begin{column}{0.48\textwidth}
			\setbeamercolor{block title example}{fg=white, bg=gray!70!black}
			\begin{exampleblock}{\small \textbf{STABILITY BREAKTHROUGHS}}
				\vspace{0.5ex}
				\small
				\textbullet~\textbf{2025:} Stability for active-circuit ROMs.\\
				\textbullet~\textbf{2017:} Non-linear network foundations.\\
				\textcolor{red!80!black}{\textbf{Result:}} Eliminated the "unstable model" industrial bottleneck.
			\end{exampleblock}
		\end{column}

	\end{columns}

	\vfill
	\centering
	\textcolor{gray!50}{\rule{0.6\textwidth}{0.4pt}} \\
	\small \textcolor{darkgray}{Research validated by over 15 years of peer-reviewed excellence.}
\end{frame}

!black}
% 	\setbeamercolor{block body example}{bg=green!5, fg=black}
% 	\begin{exampleblock}{\centering \scriptsize \textbf{KEY BREAKTHROUGHS}}
% 		\centering
% 		\vspace{1ex}
% 		{\huge \textbf{60$\times$}} \\
% 		\footnotesize \textbf{Total Speedup} \\
% 		\vspace{1.5ex}
% 		\textcolor{green!30}{\rule{0.8\textwidth}{0.5pt}} \\
% 		\vspace{1.5ex}
% 		\scriptsize
% 		\textbf{Guaranteed Preservation:} \\
% 		\textcolor{blue!80!black}{\textbf{Asymptotic Stability}} \\
% 		\textit{(By Construction)}

% 		\vspace{2ex}

% 		\textcolor{blue!80!black}{\textbf{Numerical Robustness}} \\
% 		\textit{Well-Conditioned ROMs}
% 		\vspace{1ex}
% 	\end{exampleblock}
% 	\end{column}
% 	\end{columns}

% 	\vfill

% 	%% Visionary Footer
% 	\begin{beamercolorbox}[wd=\textwidth, sep=1.2ex, center, rounded=true]{block body example}
% 		\large \itshape \textcolor{red!80!black}{"Moving beyond 'fast' simulation to 'guaranteed' simulation for large-scale active networks."}
% 	\end{beamercolorbox}
% 	\vspace{1ex}
% \end{frame}

% %%%%%%%%%%%%%%%%%%%%%%%%%%%%%%%%%%%%%%%%%%%%%%%%%%%%%%%%%%%%%%%%%%
% %%% Slide: EM-Circuit Co-Simulation (Delayed PEEC Systems)
% %%%%%%%%%%%%%%%%%%%%%%%%%%%%%%%%%%%%%%%%%%%%%%%%%%%%%%%%%%%%%%%%%%
% \begin{frame}[t]{EM-Circuit Co-Simulation: Delayed PEEC Systems}
% 	\framesubtitle{\textcolor{red}{PMOR and NILT for High-Frequency Retarded Networks}}

% 	\lmv{-1}
% 	\begin{mybox}
% 		B. Nouri, L. Lombardi, et al., ``Parameterized model order reduction of delayed PEEC circuits,''
% 		\textit{IEEE Transactions on Electromagnetic Compatibility}, vol. 62, no. 3, Jun. 2020.
% 	\end{mybox}

% 	\vspace{1ex}

% 	\begin{columns}[T, onlytextwidth]

% 		%% LEFT COLUMN: Technical Narrative
% 		\begin{column}{0.62\textwidth}
% 			\small
% 			\textbf{\textcolor{blue!80!black}{The Challenge: Neutral Delay DEs (NDDEs)}}
% 			\begin{itemize}
% 				\item High-speed retardation requires Partial Element Equivalent Circuits.
% 				\item Time-stepping is bottlenecked by intensive interpolation.
% 				\item Traditional NDDE solvers suffer from passivity/stability crises.
% 			\end{itemize}

% 			\vspace{2ex}

% 			\textbf{\textcolor{blue!80!black}{The Innovation: Decoupled PMOR + NILT}}
% 			\begin{itemize}
% 				\item \textbf{Compact Subspaces:} Frequency-independent moment matching reduces ROM size to minimal dimensions.
% 				\item \textbf{NILT Solver:} Direct Laplace-domain evaluation, eliminating time-domain interpolation overhead entirely.
% 			\end{itemize}
% 		\end{column}

% 		%% RIGHT COLUMN: The "Hero" Metric
% 		\begin{column}{0.34\textwidth}
% 			\setbeamercolor{block title alert}{bg=red!20!black, fg=white}
% 			\begin{alertblock}{\centering \textbf{MEASURED IMPACT}}
% 				\centering
% 				\vspace{1ex}
% 				{\huge \textbf{1000$\times$}} \\
% 				\footnotesize \textbf{Simulation Speedup} \\
% 				\vspace{1.5ex}
% 				\textcolor{gray!50}{\rule{0.8\textwidth}{0.4pt}} \\
% 				\vspace{1.5ex}
% 				\small \textbf{NILT vs. NDDE} \\
% 				\footnotesize Robust Solver Execution
% 				\vspace{1ex}
% 			\end{alertblock}
% 		\end{column}
% 	\end{columns}

% 	\vfill

% 	%% Visionary Footer
% 	\begin{beamercolorbox}[wd=\textwidth, sep=1.2ex, center, rounded=true]{block body example}
% 		\large \itshape \textcolor{red!80!black}{"Scaling retarded system simulation by three orders of magnitude for industrial SI/PI verification."}
% 	\end{beamercolorbox}
% 	\vspace{1ex}
% 	\vfill
% \end{frame}

%%%%%%%%%%%%%%%%%%%%%%%%%%%%%%%%%%%%%%%%%%%%%%%%%%%%%%%%%%%%%%%%%%
%%% SLIDE 1: UQ - Nonlinear Systems
%%%%%%%%%%%%%%%%%%%%%%%%%%%%%%%%%%%%%%%%%%%%%%%%%%%%%%%%%%%%%%%%%%
\begin{frame}[t]{UQ: Efficient Simulation of Nonlinear Systems}
	\vspace{1ex}

	% Consistent Citation Style
	\begin{beamercolorbox}[wd=\textwidth, sep=1ex, rounded=true]{block body}
		\scriptsize B.~Nouri, E.~Gad, A.~Nouri, and M.~Nakhla, ``DC-centric parameterized reduced-order model via moment-based interpolation projection (MIP) algorithm,'' \emph{IEEE Trans. CPMT}, vol.~10, no.~8, 2020.
	\end{beamercolorbox}

	\vspace{2ex}

	\begin{columns}[T, onlytextwidth]
		\begin{column}{0.62\textwidth}
			\small
			\textbf{\textcolor{blue!80!black}{Challenge: The Newton-Raphson Bottleneck}}
			\begin{itemize}
				\item Standard Monte Carlo is too slow for SI/PI \& RF.
				\item Iterative solvers recompute DC points for every sample.
				\item Evaluation of reduced nonlinear terms is the primary cost.
			\end{itemize}

			\vspace{2ex}
			\textbf{\textcolor{blue!80!black}{Solution: Parameterized DEIM Formulation}}
			\begin{itemize}
				\item Developed a surrogate using \textbf{DC-centric DEIM}.
				\item Directly exploits structurally parameterized moments.
			\end{itemize}
		\end{column}

		\begin{column}{0.34\textwidth}
			\setbeamercolor{block title alert}{bg=red!70!black, fg=white}
			\setbeamercolor{block body alert}{bg=red!5, fg=black}
			\begin{alertblock}{\centering \textbf{MEASURED IMPACT}}
				\centering
				\vspace{1ex}
				{\huge \textbf{50$\times$}} \\
				\footnotesize \textbf{Simulation Speed-up} \\
				\vspace{1.5ex}
				\textcolor{red!20}{\rule{0.8\textwidth}{0.5pt}} \\
				\vspace{1.5ex}
				\small \textbf{Zero-Risk} \\
				\footnotesize Tape-out Readiness
				\vspace{1ex}
			\end{alertblock}
		\end{column}
	\end{columns}

	\vfill
	\begin{beamercolorbox}[wd=\textwidth, sep=1.2ex, center, rounded=true]{block body}
		\small \itshape \textcolor{blue!80!black}{"Moving UQ from a 'prohibitive cost' to a 'practical design-stage tool' for SI/PI \& RF systems."}
	\end{beamercolorbox}
\end{frame}

%%%%%%%%%%%%%%%%%%%%%%%%%%%%%%%%%%%%%%%%%%%%%%%%%%%%%%%%%%%%%%%%%%
%%% SLIDE 2: HB - RFIC Innovation
%%%%%%%%%%%%%%%%%%%%%%%%%%%%%%%%%%%%%%%%%%%%%%%%%%%%%%%%%%%%%%%%%%
\begin{frame}[t]{HB: Efficient Periodic Steady-State (PSS) Analysis}
	\vspace{1ex}

	\begin{beamercolorbox}[wd=\textwidth, sep=1ex, rounded=true]{block body}
		\scriptsize B.~Nouri, E.~Gad, A.~Nouri, and M.~Nakhla, ``Parameterized periodic steady-state analysis via reduced-order modelling,'' \emph{IEEE Trans. MTT}, vol.~69, no.~11, 2021.
	\end{beamercolorbox}

	\vspace{2ex}

	\begin{columns}[T, onlytextwidth]
		\begin{column}{0.62\textwidth}
			\small
			\textbf{\textcolor{blue!80!black}{Problem: The "Transformation Tax"}}
			\begin{itemize}
				\item Traditional solvers "ping-pong" between Time and Frequency domains for nonlinear terms.
				\item FFT/IFFT operations dominate the CPU time.
			\end{itemize}

			\vspace{2ex}
			\textbf{\textcolor{blue!80!black}{Innovation: Direct Harmonic Domain DEIM}}
			\begin{itemize}
				\item Evaluation of surrogates directly in the harmonic domain.
				\item Eliminates domain conversion overhead entirely.
				\item Successfully benchmarked on large-scale RF designs.
			\end{itemize}
		\end{column}

		\begin{column}{0.34\textwidth}
			\setbeamercolor{block title alert}{bg=red!70!black, fg=white}
			\setbeamercolor{block body alert}{bg=red!5, fg=black}
			\begin{alertblock}{\centering \textbf{MEASURED IMPACT}}
				\centering
				\vspace{1ex}
				{\huge \textbf{30$\times$}} \\
				\footnotesize \textbf{Total Speedup} \\
				\vspace{1.5ex}
				\textcolor{red!20}{\rule{0.8\textwidth}{0.5pt}} \\
				\vspace{1.5ex}
				\small \textbf{Days $\rightarrow$ Hours} \\
				\footnotesize Design Cycle Reduction
				\vspace{1ex}
			\end{alertblock}
		\end{column}
	\end{columns}

	\vfill
	\begin{beamercolorbox}[wd=\textwidth, sep=1.2ex, center, rounded=true]{block body}
		\small \itshape \textcolor{blue!80!black}{"Transforming infeasible multi-day simulations into routine, automated verification tasks for high-yield RFICs."}
	\end{beamercolorbox}
\end{frame}

%%%%%%%%%%%%%%%%%%%%%%%%%%%%%%%%%%%%%%%%%%%%%%%%%%%%%%%%%%%%%%%%%%
%%% SLIDE 3: Stability - Active Circuits
%%%%%%%%%%%%%%%%%%%%%%%%%%%%%%%%%%%%%%%%%%%%%%%%%%%%%%%%%%%%%%%%%%
\begin{frame}[t]{Stability by Construction: Active Circuit Macromodeling}
	\vspace{1ex}

	\begin{beamercolorbox}[wd=\textwidth, sep=1ex, rounded=true]{block body}
		\scriptsize G. Ghaly, E. Gad, M. Nakhla, and \textbf{B. Nouri}, ``An Algorithmic Approach to Formulate Well-Conditioned Stable Reduced-Order Models of Active Circuits,'' \textit{IEEE Trans. CPMT}.
	\end{beamercolorbox}

	\vspace{2ex}

	\begin{columns}[T, onlytextwidth]
		\begin{column}{0.62\textwidth}
			\small
			\textbf{\textcolor{blue!80!black}{Problem: Simulation Divergence}}
			\begin{itemize}
				\item Active circuits yield unstable ROMs under Galerkin projection.
				\item Results in catastrophic failure during transient analysis.
			\end{itemize}

			\vspace{2ex}
			\textbf{\textcolor{blue!80!black}{Innovation: Lyapunov-Based Projection}}
			\begin{itemize}
				\item Left-projection ($V$) satisfies \textbf{Lyapunov stability}.
				\item Efficient \textbf{generalized eigenproblem} avoids optimization.
				\item Built-in correction for numerical robustness.
			\end{itemize}
		\end{column}

		\begin{column}{0.34\textwidth}
			\setbeamercolor{block title alert}{bg=red!70!black, fg=white}
			\setbeamercolor{block body alert}{bg=red!5, fg=black}
			\begin{alertblock}{\centering \textbf{MEASURED IMPACT}}
				\centering
				\vspace{1ex}
				{\huge \textbf{60$\times$}} \\
				\footnotesize \textbf{Total Speedup} \\
				\vspace{1.5ex}
				\textcolor{red!20}{\rule{0.8\textwidth}{0.5pt}} \\
				\vspace{1.5ex}
				\small \textbf{Guaranteed Stability} \\
				\footnotesize Numerical Robustness
				\vspace{1ex}
			\end{alertblock}
		\end{column}
	\end{columns}

	\vfill
	\begin{beamercolorbox}[wd=\textwidth, sep=1.2ex, center, rounded=true]{block body}
		\small \itshape \textcolor{blue!80!black}{"Moving beyond 'fast' simulation to 'guaranteed' simulation for large-scale active networks."}
	\end{beamercolorbox}
\end{frame}

%%%%%%%%%%%%%%%%%%%%%%%%%%%%%%%%%%%%%%%%%%%%%%%%%%%%%%%%%%%%%%%%%%
%%% SLIDE 4: EM-Circuit (Delayed PEEC)
%%%%%%%%%%%%%%%%%%%%%%%%%%%%%%%%%%%%%%%%%%%%%%%%%%%%%%%%%%%%%%%%%%
\begin{frame}[t]{EM-Circuit Co-Simulation: Delayed PEEC Systems}
	\vspace{1ex}

	\begin{beamercolorbox}[wd=\textwidth, sep=1ex, rounded=true]{block body}
		\scriptsize B. Nouri, L. Lombardi, et al., ``Parameterized model order reduction of delayed PEEC circuits,'' \emph{IEEE Trans. EMC}, vol.~62, no.~3, 2020.
	\end{beamercolorbox}

	\vspace{2ex}

	\begin{columns}[T, onlytextwidth]
		\begin{column}{0.62\textwidth}
			\small
			\textbf{\textcolor{blue!80!black}{Challenge: Neutral Delay DEs (NDDEs)}}
			\begin{itemize}
				\item High-speed retardation requires complex PEEC circuits.
				\item Time-stepping is bottlenecked by interpolation.
				\item Traditional solvers suffer from passivity crises.
			\end{itemize}

			\vspace{2ex}
			\textbf{\textcolor{blue!80!black}{Innovation: Decoupled PMOR + NILT}}
			\begin{itemize}
				\item \textbf{Compact Subspaces:} Frequency-independent matching.
				\item \textbf{NILT Solver:} Direct Laplace-domain evaluation, eliminating time-domain interpolation overhead.
			\end{itemize}
		\end{column}

		\begin{column}{0.34\textwidth}
			\setbeamercolor{block title alert}{bg=red!70!black, fg=white}
			\setbeamercolor{block body alert}{bg=red!5, fg=black}
			\begin{alertblock}{\centering \textbf{MEASURED IMPACT}}
				\centering
				\vspace{1ex}
				{\huge \textbf{1000$\times$}} \\
				\footnotesize \textbf{Simulation Speedup} \\
				\vspace{1.5ex}
				\textcolor{red!20}{\rule{0.8\textwidth}{0.5pt}} \\
				\vspace{1.5ex}
				\small \textbf{NILT vs. NDDE} \\
				\footnotesize Three-Order Scale Shift
				\vspace{1ex}
			\end{alertblock}
		\end{column}
	\end{columns}

	\vfill
	\begin{beamercolorbox}[wd=\textwidth, sep=1.2ex, center, rounded=true]{block body}
		\small \itshape \textcolor{blue!80!black}{"Scaling retarded system simulation by three orders of magnitude for industrial SI/PI verification."}
	\end{beamercolorbox}
\end{frame}
